% Evans 12.4, printed pp.636--642; English PDF pp.653--659.
% Chinese start: printed p.451 / scan p.469.
\section{Critical Power Nonlinearity}
\label{sec:12-4}
\subsection{Scaling and the external linear estimate}
In $\R^3$ the defocusing equation
\begin{equation}\label{eq:12-quintic}
 \PDEWave u+u^5=0,\qquad
 E=\int\{(u_t^2+|Du|^2)/2+u^6/6\}
\end{equation}
is energy critical. For the general power,
$u_\lambda(x,t)=\lambda^{2/(p-1)}u(\lambda x,\lambda t)$ solves
the same equation, while its energy scales by
$\lambda^{4/(p-1)-1}$. At $p=5$ this factor is one: a small spatial
scale cannot be excluded by total energy alone.
\PDEMarginDiagram{12.M}{能量临界的缩放指数}
{图中 $a(p)=4/(p-1)-1$，且 $E[u_\lambda]=\lambda^{a(p)}E[u]$。
$p=5$ 是三维退聚焦波方程的能量临界指数；
它与后面聚焦小数据爆破的 $1+\sqrt2$ 阈值属于不同问题。}
{ch12-scaling-margin}

\begin{theorem}[Linear estimate; source \S12.4, Theorem 1]
For $\PDEWave v=G$ on a time interval $I=[s,\tau]$ in $\R^3$,
assume $v(s)\in\dot H^1$, $v_t(s)\in L^2$, and
$G\in L^1(I;L^2)$. Then
\begin{equation}\label{eq:12-strichartz}
 \begin{split}
 &\|Dv,v_t\|_{L^\infty(I;L^2)}+\|v\|_{L^\infty(I;L^6)}
       +\|v\|_{L^4(I;L^{12})}\\
 &\hspace{4mm}\le C\{\|Dv(s)\|_2+\|v_t(s)\|_2
                               +\|G\|_{L^1(I;L^2)}\}.
 \end{split}
\end{equation}
We explicitly use the homogeneous three-dimensional Strichartz
estimate for the $(4,12)$ pair as an \emph{external theorem}, with
constant independent of the interval. It is not proved here.
\end{theorem}
\begin{expandedproof}[the proved energy part and the external input]
The energy argument of \S\ref{sec:12-2}, without the $L^2$ position
term, gives the first norm on the left. The homogeneous Sobolev
inequality $\|v\|_6\le C\|Dv\|_2$ gives the second. The external
homogeneous estimate bounds the third norm for free waves by their
$\dot H^1\times L^2$ data. Apply this to the free evolution starting
at each forcing time and use Minkowski in Duhamel's integral. Its
$\dot H^1\times L^2$ size is $\|G(t)\|_2$, giving the forcing term
in \eqref{eq:12-strichartz}. Restriction of the homogeneous estimate
on the whole time line explains uniformity under interval shortening.
The dispersive proof of that homogeneous estimate remains external;
energy conservation alone does not imply an $L^4_tL^{12}_x$ bound.
\end{expandedproof}

\subsection{Morawetz identity and potential nonconcentration}
Fix an apex $(x_0,T)$ before which $u$ is smooth and has finite
conserved energy. Set $x_0=0$ by translation and $r=|x|$.
The local energy $E_K(t)$ decreases, so it has a limit $E_*\ge0$
as $t\uparrow T$. Define the tail flux
\begin{equation}\label{eq:12-tail-flux}
 \phi(s)=E_K(s)-E_*
 =\frac1{\sqrt2}\int_{\Gamma_{s,T}}
       \{\tfrac12|u_t\nu-Du|^2+u^6/6\}\,dS.
\end{equation}
Monotone convergence in the truncated flux identity gives this
formula and $\phi(s)\to0$. At this stage $E_*=0$ is \emph{not} known.

\begin{lemma}[Multiplier identity; source equations (6)--(8)]
Put $\theta=t-T$, $m=\theta u_t+x\cdot Du+u$ and
\[
 P=\theta\{(u_t^2+|Du|^2)/2+u^6/6\}+(x\cdot Du+u)u_t,
\]
\[
 Q=mDu+x\{(u_t^2-|Du|^2)/2-u^6/6\}.
\]
Then $P_t-\operatorname{div}Q=m(\PDEWave u+u^5)-u^6/3$.
\end{lemma}
\PDEAnnotation{先证明通量尾部趋零}{正确的关系是
$\phi(s)=E_K(s)-\lim_{t\uparrow T}E_K(t)$。
单调性保证尾部趋零，不要求顶端剩余能量预先为零。
后面通过线性 Duhamel 的强能量极限再证明全部能量不集中。}
\begin{expandedproof}[an identity also valid off solutions]
Use $\theta_t=1$, $\operatorname{div}x=3$ and
$D(x\cdot Du)=Du+D^2u\,x$.
The mixed terms $\theta Du\cdot Du_t$ cancel, as do
$x\cdot D(u_t)u_t$ and $x\cdot D^2u\,Du$ after differentiation.
The remaining second derivatives form $m(u_{tt}-\Delta u)$.
The potential terms are $m u^5-u^6/3$: their residual coefficient
is $1/6+3/6-1=-1/3$. The coefficients of $u_t^2$ and $|Du|^2$
are both zero. This proves the identity without presupposing the PDE.
\end{expandedproof}

\begin{theorem}[Potential nonconcentration; source \S12.4, Theorem 2]
For $R=T-s$,
\begin{equation}\label{eq:12-nonconcentration}
 \int_{B(x_0,R)}|u(x,s)|^6dx
       \le C\{\phi(s)+\phi(s)^{1/3}\}\longrightarrow0.
\end{equation}
Since $\phi\le E$, the right side can also be written
$C_E\phi(s)^{1/3}$, making its energy dependence explicit.
\end{theorem}
\begin{expandedproof}[top term, boundary cancellation, and the endpoint]
Translate $s$ to zero, so $R$ replaces $T$. Integrate the multiplier
identity on $K_{0,\tau}$ and let $\tau\uparrow R$.
The top integral tends to zero even if the apex is singular:
with $\rho=R-\tau$,
\[
 \left|\int_{B(0,\rho)}P(x,\tau)dx\right|
 \le C\rho E+C\|u\|_{L^2(B_\rho)}\|u_t\|_2
 \le C\rho(E+E^{2/3}).
\]
Here Holder gives $\|u\|_{L^2(B_\rho)}\le C\rho\|u\|_6$.
Thus, putting
$A=\frac1{\sqrt2}\int_\Gamma(P-Q\cdot\nu)dS$ and
$B=\int_{B_R}P(x,0)dx$, we obtain $A-B=-\iint_Ku^6/3\le0$.

On the mantle $r=R-t$,
$P-Q\cdot\nu=-r(u_t-u_r)^2+u(u_t-u_r)$.
Set $w(y)=u(y,R-|y|)$. Then $w_r=u_r-u_t$ and
\begin{equation}\label{eq:12-mantle-square}
 A=-\int_{B_R}r\left(w_r+\frac wr\right)^2dy
                      +\frac12\int_{\partial B_R}u(x,0)^2dS_x.
\end{equation}
Indeed $A=-\int(rw_r^2+ww_r)$ and
$\int(w^2/r+ww_r)=\frac12\int_{\partial B_R}w^2$, from
$\operatorname{div}(y/r)=2/r$.
This calculation is valid for the possibly singular cone trace:
flux gives $Dw\in L^2$ and $w\in L^6$, so $w\in H^1(B_R)$
(the point at zero is removable for $H^1$ in dimension three).
The weighted term $w^2/r$ is integrable by Holder. Approximate in
$H^1$ with the outer smooth trace fixed, or use radial cutoffs and
a vanishing inner-boundary subsequence, to justify the integration.

For the bottom term use Young's inequality and $r\le R$:
\[
 \begin{split}
 B&\le-\frac R6\int_{B_R}u^6
       +\frac R2\int_{B_R}(|Du+xu/r^2|^2-|Du|^2)\\
  &=-\frac R6\int_{B_R}u^6+\frac12\int_{\partial B_R}u^2.
 \end{split}
\]
The last equality follows from $\operatorname{div}(x/r^2)=1/r^2$;
the origin contributes no boundary term for the smooth bottom data.
The boundary terms cancel between this bound and
\eqref{eq:12-mantle-square}. Therefore
\[
 R\int_{B_R}u(x,0)^6dx
 \le C\int_{B_R}r(w_r+w/r)^2
 \le CR\phi(0)+C\int_{B_R}w^2/r.
\]
Finally
\[
 \int_{B_R}w^2/r
 \le\left(\int_{B_R}w^6\right)^{1/3}
      \left(\int_{B_R}r^{-3/2}\right)^{2/3}
 \le CR\phi(0)^{1/3}.
\]
Divide by $R$ and undo the time translation. Every estimate is on
truncated smooth cones before a justified limit, so no regularity
at the candidate blow-up point has been assumed.
\end{expandedproof}

\subsection{Localizing Strichartz and closing the global proof}
\begin{lemma}[Localized estimate and the bootstrap barrier]
Let $C=\{s<t<\tau,\ |x-x_0|<R-(t-s)\}$ be a backward cone.
For a smooth quintic solution set
$X(\tau)=\|u\|_{L^4(s,\tau;L^{12}(C_t))}$.
Then
\begin{equation}\label{eq:12-cone-strichartz}
 X(\tau)\le C_0+C\sup_{s<t<\tau}\|u(t)\|_{L^6(C_t)}X(\tau)^4,
 \qquad C_0\le C\sqrt E.
\end{equation}
If the displayed $L^6$ coefficient is sufficiently small depending
only on $E$, then $X(\tau)$ stays bounded up to the cone apex.
\end{lemma}
\begin{expandedproof}[no unsupported restriction of a global estimate]
Solve the linear equation with data $(u(s),u_t(s))$ on all space
and forcing $-\mathbf1_Cu^5$, zero outside $C$. Finite propagation
for the linear equation implies its solution equals $u$ inside $C$:
the backward cone of any point of $C$ remains inside $C$ down to $s$.
Apply \eqref{eq:12-strichartz} to this linear solution. Holder in
space and time gives
\[
 \|u^5\|_{L^1L^2(C)}
 \le\sup_t\|u(t)\|_{L^6(C_t)}
                 \int_s^\tau\|u(t)\|_{L^{12}(C_t)}^4dt.
\]
This proves the estimate using a globally defined forcing and data;
no boundary Strichartz theorem is being assumed.
Choose $B=2\max(C_0,1)$ and require the nonlinear coefficient to be
less than $1/(4B^3)$. The continuous function $X$ starts at zero.
At a first crossing of $B$, the inequality would give
$B\le C_0+B/4\le3B/4$, impossible. It therefore remains below $B$.
The small coefficient alone is insufficient unless this continuity
argument excludes the large branch of $X\le C_0+\varepsilon X^4$.
\end{expandedproof}

\begin{theorem}[Global smooth quintic solution; source \S12.4, Theorem 3]
For $g,h\in C_c^\infty(\R^3)$, the defocusing quintic problem has
a unique global smooth solution. The proof uses the stated external
homogeneous Strichartz estimate.
\end{theorem}
\begin{expandedproof}[from potential flux to an actual continuation bound]
Assume a finite maximal time $T$. For each $x_0$, potential
nonconcentration makes the $L^6$ coefficient in
\eqref{eq:12-cone-strichartz} small on a sufficiently short tail of
the cone with apex $(x_0,T)$. Thus $X(T)<\infty$ on that cone, and
$G=-\mathbf1_Ku^5\in L^1(s,T;L^2)$.
The linear solution with this forcing and full data at $s$ agrees
with $u$ inside the cone, as in the lemma. Because its forcing is
integrable in time, Duhamel gives a strong
$\dot H^1\times L^2$ limit at $T$.
Strong $L^2$ convergence of its gradient and velocity, together with
absolute continuity of their squared integrals at the limiting time,
now implies
\[
 \int_{B(x_0,T-t)}(|Du|^2+u_t^2)dx\longrightarrow0.
\]
Together with \eqref{eq:12-nonconcentration} this proves
$E_K(t)\to0$. This argument replaces an unproved differentiation
estimate on a cone with a characteristic boundary.

Fix $x_0$. Choose $s<T$ so its local energy is smaller than a chosen
$\varepsilon>0$. At this fixed smooth time, absolute continuity in
the radius allows a $\delta>0$ such that the energy on
$B(x_0,T-s+\delta)$ is less than $2\varepsilon$.
Positive cone flux propagates this bound into the enlarged cone
$|x-x_0|<T-t+\delta$. Its $L^6$ norm is at most
$(12\varepsilon)^{1/6}$. The localized lemma, still using the full
initial energy $E$ for $C_0$, gives a finite $L^4L^{12}$ norm on this
enlarged cone. Choose $\varepsilon$ small enough for the fixed barrier.

Finite propagation bounds all supports for $t<T$ by the fixed closed
ball $\overline B(0,R_0+T)$. Cover that ball by finitely many of the
balls $B(x_j,\delta_j/2)$ just constructed. These balls lie inside
the corresponding enlarged sections for every $t<T$. Start after
the largest of their finitely many start times, and use the triangle
inequality in $L^{12}$ and then $L^4$. Adding the earlier smooth
strip proves $u\in L^4(0,T;L^{12}(\R^3))$.
No object $\operatorname{spt}u(\cdot,T)$ is needed before existence
at $T$ is known.

Finally $v=D_ju$ solves $\PDEWave v=-5u^4v$. The homogeneous energy
estimate and Sobolev give, for
$A_j(t)=\|Dv(t)\|_2+\|v_t(t)\|_2$,
\[
 A_j(t)\le C_j+C\int_0^t\|u(s)\|_{12}^4 A_j(s)ds,
 \qquad \|u^4v\|_2\le\|u\|_{12}^4\|v\|_6.
\]
Gronwall bounds $A_j$ up to $T$. Thus $Du$ is bounded in $L^6$;
energy already bounds $u$ in $L^6$. The embedding
$W^{1,6}(\R^3)\hookrightarrow L^\infty$ supplies a uniform amplitude
bound. The continuation theorem of \S\ref{sec:12-3} contradicts
maximality. Local uniqueness extends to all time.
\end{expandedproof}
\PDEBodyDiagram{12.3}{Why the cone must be enlarged}
{The flux tail first makes the potential energy small on shrinking cones.
Localized Strichartz and a linear Duhamel limit then exclude concentration
of the kinetic and gradient energy. At a fixed time $s$, the small-energy
ball can be enlarged by $\delta>0$. The enlarged sections retain a ball
of radius $\delta$ near $T$, allowing a finite cover of the fixed possible
support. A cone with zero terminal radius alone would not provide that cover.}
{ch12-nonconcentration}
